\documentclass[11pt]{article}
\usepackage[margin=1in]{geometry}
\usepackage{booktabs,longtable,hyperref,amsmath}
\usepackage[T1]{fontenc}
\hypersetup{colorlinks=true,linkcolor=blue,urlcolor=blue}
\title{\textbf{Response to the review of Draft 23}}
\author{}
\date{11 October 2026}
\begin{document}
\maketitle
\vspace{-3em}

Five items. Three applied, one applied in a different form than proposed because the
proposed form rests on a claim the manuscript got wrong, one not applicable to me.
All three checkers green afterwards: 63 headline quantities, 301 printed values,
25 superseded literals.

\section*{1. Hull-distance matching --- applied as a correction, not as the computation}

The computation was not run. The reason is that the manuscript's own account of what it
would take was wrong, and that is the part worth fixing first.

Draft 23 said the sub-hull margin ``is recoverable by removing each entry from its own
phase diagram and recomputing the hull \ldots\ That is a few hours of computation on
entries we already hold.'' We do not hold those entries. The cached Materials Project pull
is four columns --- \texttt{material\_id}, \texttt{energy\_above\_hull},
\texttt{nelements}, \texttt{chemsys} --- for 152{,}427 GGA entries and 37{,}124 r$^2$SCAN
entries. Reconstructing a phase diagram needs the competing entries of each chemical
system, which is a different and much larger retrieval. The ``few hours'' figure was an
estimate of a cost I had not measured, for an operation on data I did not have.

I also could not run the retrieval from here: \texttt{api.materialsproject.org} is
unreachable from both environments available to this session (403 at the proxy from the
Linux VM, connection refused from the container), and neither has \texttt{mp-api} or
\texttt{pymatgen} installed. The pull has to happen on the machine that ran the original
queries.

So the passage has been rewritten to say what is true and nothing more: the comparison is
not matched on hull depth; the quantity exists on the GNoME side and not on the Materials
Project side as retrieved, because \texttt{energy\_above\_hull} is floored at zero for
stable entries; recovering it requires the competing entries of each system; it was not
done; and no cost estimate is offered. The ``first thing we would add'' sentence is gone,
which removes the visibly-unfinished reading.

The downgrade asked for is also in, as a single new limitation paragraph --- distance below
the hull predicts disagreement within GNoME, the two sets are not matched on it, this is
the largest unaddressed confound in that comparison, and it is why the comparison is
exploratory. Stated once in \S4.13 and once, briefly, in Limitations.

\section*{2. ``All'' reference headline cells --- computed}

This was the right request and it was cheap, because the ``All'' column is in the released
file. New \texttt{scripts/phase40\_allref\_cells.py} recomputes the two cells the abstract
quotes, with one prototype draw per replicate shared across both cells so that the spread
and ratio intervals account for prototypes appearing in both. 20{,}000 replicates. It
reproduces the global ``All'' rate of 25.8\% as a sanity check.

\begin{center}
\begin{tabular}{lrrr}
\toprule
 & $n$ & MP reference & ``All'' reference \\
\midrule
oxide, no Ln, 3$d$ (4-el)        & 2{,}478 & 30.7\% & 37.4\% [34.6, 40.4] \\
intermetallic, no Ln/3$d$ (4-el) & 5{,}025 & 4.0\%  & 14.5\% [12.9, 16.1] \\
\midrule
spread (pp)                      &         & $+26.7$ [$+23.5$, $+29.9$] & $+22.9$ [$+19.7$, $+26.3$] \\
ratio                            &         & 7.67 [6.29, 9.52] & 2.58 [2.26, 2.97] \\
\bottomrule
\end{tabular}
\end{center}

\textbf{The headline survives.} Both rates rise, the ratio compresses from 7.67-fold to
2.58-fold, but the absolute spread falls only from $+26.7$ to $+22.9$ points and its
interval [$+19.7$, $+26.3$] is nowhere near zero.

Worth noting against the review's framing: the cell-level compression (7.67 to 2.58) is
much milder than the anion-class compression of Table~13 (2.1 to 1.3). The class figures
pool over element count and over lanthanide and 3$d$ content, and the intermetallic class
as a whole gains more from the more inclusive reference set than its lanthanide- and
3$d$-free quaternary subset does. So the class-level number overstates how much the
headline moves. That is now said in the text.

One thing the new table deliberately does not do is print a second interval for the
MP-reference cells. The shared-draw design gives $[27.6, 33.8]$ for the oxide cell against
the canonical per-cell $[27.7, 33.8]$, and \texttt{check\_stale.py} refused the build
because 27.6 is a pinned superseded value from the pre-convergence 600-replicate runs. It
was right to refuse: two intervals for the same quantity differing in the last digit carry
no information and invite exactly that confusion. The MP column is point estimates only,
with the caption saying why.

\section*{3. Ni intermetallic interval --- fixed, and the inference about the checker was right}

Confirmed and corrected. The row for Ni in intermetallics in
\texttt{within\_anion\_excess\_ci.csv} reads $n=4979$, excess $2.0$, interval $0.3$ to
$3.7$: the table was right, the prose was wrong, and it now reads $+2.0$ [$+0.3$, $+3.7$].

The second half of the observation was the more useful one. \texttt{check\_paper.py}
filtered that file to iron rows only, so no nickel or cobalt
endpoint was ever compared, and the value was in prose, which the table parser cannot see
and which \texttt{check\_stale.py} only covers for values already known to be superseded.
Three things changed:

\begin{itemize}
\item \texttt{check\_paper.py} now carries a \texttt{PROSE\_INTERVALS} list: each interval
quoted outside a table is matched by regex and diffed against the CSV. The Ni sentence is
the first entry. Verified by regression --- restoring 3.8 produces a reported mismatch
(\mbox{paper 3.8}, \mbox{results 3.7}) and exit status 1.
\item \texttt{check\_stale.py} pins the wrong interval string at zero permitted
occurrences.
\item The new two-reference table is parsed too, so its 12 values cannot go stale.
\end{itemize}

Printed-value coverage went from 278 to 301. \texttt{VALIDATION.md} now states plainly
that prose intervals are enumerated rather than found automatically, and that a new one
needs a line added.

I had also told the author that \texttt{check\_paper.py} covered ``the endpoints of every
interval in the within-anion excess table.'' It covered Fe's. That overclaim is corrected
in \texttt{VALIDATION.md}.

\section*{4. External review --- not mine to apply}

Correct and unchanged. The manuscript says the core definitions and scripts have not been
reviewed by anyone other than the author, and that stays true until someone reviews them.
\texttt{REVIEWER.md} already orders the eight areas to inspect, and the seven items listed
in the review map onto B1--B7 of it, so the ask can be forwarded as-is.

\section*{5. Residual process narration --- applied}

Removed from the manuscript and left in the repository:

\begin{itemize}
\item the reference-run taxonomy with its counts and the ``expected to be updated whenever
a number legitimately changes'' explanation;
\item the per-layer description of what each checker parses;
\item ``cheaper than the old implementation's 400,'' now stated as a property of the
aggregation rather than as a comparison against a previous version;
\item ``retains the \ldots\ superseded independence-based intervals,'' now stated as what
\texttt{phase28} carries and what \S4.13 quotes it for;
\item ``records the defects found in preparing this work as calibration.''
\end{itemize}

What is kept is one short paragraph: the driver regenerates and checks, CI runs it on a
clean runner after deleting the stored results, \texttt{VALIDATION.md} has the detail, and
reproducing a number is not the same as the number being right. That last sentence is a
warning about how to read the paper, not software maintenance, so it belongs in the paper.

\section*{Not applied}

Nothing was rejected outright. The only departure from the review is item 1, where the
computation was replaced by a correction of the manuscript's claim about the computation,
for the reasons given. If the hull-depth retrieval is run later, the matched comparison
should be redone on continuous depth rather than on bands, and the result will either
strengthen \S4.13 or weaken it --- the section is already labelled exploratory either way.

\end{document}
